% notes/camacho_2025.tex
% Fiche de lecture — Camacho de la Rosa, Esquivel-Sirvent & Becerril (2025)
% Dépôt : Phonon Hydrodynamics, Hyperbolic Structure and the Spectral Inverse
%         Problem in Thin Films
% Compilation : pdflatex camacho_2025.tex  (deux passes)

\documentclass[11pt,a4paper]{article}

\usepackage[T1]{fontenc}
\usepackage[utf8]{inputenc}
\usepackage{lmodern}
\usepackage[french]{babel}
\usepackage[a4paper,top=23mm,bottom=23mm,left=24mm,right=24mm]{geometry}

\usepackage{amsmath,amssymb,mathtools}
\usepackage{booktabs,array,tabularx,longtable}
\usepackage{xcolor}
\usepackage{framed}
\usepackage{titlesec}
\usepackage{enumitem}
\usepackage{fancyhdr}
\usepackage{microtype}
\usepackage{needspace}
\usepackage{textcomp}
\usepackage{hyperref}

% ---------------------------------------------------------------- couleurs
\definecolor{filet}{gray}{0.78}
\definecolor{codebg}{HTML}{F3F4F6}
\definecolor{mut}{gray}{0.45}
\definecolor{lien}{HTML}{1F4E79}
\colorlet{shadecolor}{codebg}

\hypersetup{
  colorlinks=true, linkcolor=lien, urlcolor=lien, citecolor=lien,
  pdftitle={Camacho de la Rosa, Esquivel-Sirvent \& Becerril (2025) --- fiche de lecture},
  pdfsubject={Relaxation times of non-Fourier materials using frequency-domain thermoreflectance},
  pdfkeywords={Cattaneo-Vernotte, FDTR, temps de relaxation, sensibilite, identifiabilite}
}

% ---------------------------------------------------------------- mise en page
\setlength{\parindent}{0pt}
\setlength{\parskip}{0.55em}
\renewcommand{\arraystretch}{1.25}

\titleformat{\section}
  {\normalfont\sffamily\large\bfseries}{\thesection.}{0.55em}{}
  [\vspace{1.5pt}{\color{filet}\titlerule[0.6pt]}]
\titlespacing*{\section}{0pt}{1.6em}{0.7em}

\titleformat{\subsection}
  {\normalfont\sffamily\normalsize\bfseries}{\thesubsection}{0.55em}{}
\titlespacing*{\subsection}{0pt}{1.2em}{0.4em}

\pagestyle{fancy}
\fancyhf{}
\renewcommand{\headrulewidth}{0pt}
\renewcommand{\footrulewidth}{0pt}
\fancyfoot[L]{\footnotesize\sffamily\color{mut}Fiche de lecture \quad Camacho de la Rosa 2025 \quad passe 2}
\fancyfoot[R]{\footnotesize\sffamily\color{mut}\thepage}
\fancypagestyle{plain}{\fancyhf{}%
  \fancyfoot[L]{\footnotesize\sffamily\color{mut}Fiche de lecture \quad Camacho de la Rosa 2025 \quad passe 2}%
  \fancyfoot[R]{\footnotesize\sffamily\color{mut}\thepage}%
  \renewcommand{\headrulewidth}{0pt}}

% ---------------------------------------------------------------- macros
\newcommand{\sep}{\par\medskip\noindent{\color{filet}\rule{\linewidth}{0.6pt}}\par\medskip}
\newenvironment{encadre}{\par\smallskip\begin{shaded}\vspace{-6pt}}{\vspace{-6pt}\end{shaded}\par\smallskip}
\newcommand{\code}[1]{\texttt{\small #1}}
\newcolumntype{L}[1]{>{\raggedright\arraybackslash}p{#1}}
\newcolumntype{C}{>{\centering\arraybackslash}X}

\begin{document}

% ================================================================ titre
{\sffamily\bfseries\LARGE Camacho de la Rosa, Esquivel-Sirvent \& Becerril (2025)\\[3pt]
\large fiche de lecture\par}
\vspace{4pt}
{\color{filet}\rule{\linewidth}{1pt}}
\vspace{6pt}

\textbf{Référence.} A. Camacho de la Rosa, R. Esquivel-Sirvent, D. Becerril,
\textit{Relaxation times of non-Fourier materials using frequency-domain thermoreflectance},
Journal of Applied Physics \textbf{137}, 155103 (2025). DOI 10.1063/5.0257299.
Soumis le 9 janvier 2025, accepté le 27 mars 2025, publié le 15 avril 2025. Licence CC BY 4.0.

\textbf{Affiliations.} Centro de Investigación en Ciencias, Universidad Autónoma del Estado de
Morelos, Cuernavaca (Camacho de la Rosa)~; Instituto de Física, UNAM, Mexico
(Esquivel-Sirvent)~; Istituto di Struttura della Materia, CNR, Rome (Becerril, auteur
correspondant).

\textbf{Statut de lecture.} Passe 2, texte intégral. Matériel supplémentaire non consulté
(voir §~11). Fiche établie le 12 septembre 2026.

Article cité dans la section \textit{Prior work} du README. C'est la seule référence du corpus
qui traite explicitement de la mesure du temps de relaxation par thermoréflectance, et la seule
qui formule la sensibilité de la phase à ce paramètre. Sa lecture détermine la frontière entre
ce qui est déjà publié et ce qui reste ouvert pour la partie IV.

\sep

% ================================================================ 1
\section{Passe 1 --- les quatre réponses du protocole}

\textbf{Type d'article.} Article théorique de modélisation directe. Aucune donnée
expérimentale, aucun ajustement, aucune inversion. Le modèle de Cattaneo--Vernotte est inséré
dans le modèle direct standard d'une expérience FDTR, et la réponse en phase est calculée pour
trois configurations. Les auteurs l'énoncent eux-mêmes : «~in this work, we theoretically
investigate the viability of using frequency-domain thermoreflectance to determine non-Fourier
heat behavior~».

\textbf{Résultat en une phrase.} Pour un matériau de Cattaneo--Vernotte, la phase de la réponse
thermoréflectante acquiert une dépendance en fréquence qui disparaît sous Fourier, cette
dépendance devient sensible au voisinage de $\omega \simeq \tau^{-1}$, et le coefficient de
sensibilité $S_\tau = \mathrm{d}\phi/\mathrm{d}\ln\tau$ y prend une valeur non négligeable.

\textbf{Compatibilité des hypothèses.} Compatible sur le fond, plus restreinte sur la forme.
La loi constitutive est celle du présent travail (Cattaneo, un seul temps de relaxation, pas de
terme non local). La substitution spectrale est identique à celle codée dans
\code{src/\allowbreak forward\_model.py} (§~2.3). Les restrictions du modèle sont plus fortes : absence de
conductance d'interface, absence de pertes, absence de modèle de bruit, absence de matrice de
covariance (§~3.4).

\textbf{Faut-il continuer.} Oui, une fois, intégralement. L'article contient trois éléments qui
ne figurent nulle part ailleurs dans le corpus : la forme fermée de la phase en régime haute
fréquence (§~4.1), la définition explicite de $S_\tau$ (§~5.1), et l'aveu de bande concernant
la portée instrumentale du seuil (§~7.3).

% ================================================================ 2
\section{Rappel --- les deux modèles et la structure de la substitution}

\subsection{Équations de champ}

Les deux équations de champ utilisées sont, dans les notations de l'article (éq.~1 et 2) :
\[
\nabla^2 T(\mathbf r, t) - \frac{1}{\alpha}\,\partial_t T(\mathbf r, t)
= -\frac{1}{k}\, s_f(\mathbf r, t)
\]
\[
\nabla^2 T(\mathbf r, t) - \frac{1}{\alpha}\,\partial_t T(\mathbf r, t)
- \frac{1}{v^2}\,\partial_t^2 T(\mathbf r, t) = -\frac{1}{k}\, s_{cv}(\mathbf r, t)
\]
avec $\alpha = k/(\rho c_v)$ la diffusivité, $k$ la conductivité, et
\[
v^2 = \frac{\alpha}{\tau}
\]
la vitesse de Cattaneo--Vernotte. Les lois de flux sont $\mathbf q_f = -k\nabla T_f$ pour
Fourier et $\mathbf q + \tau\,\partial_t \mathbf q = -k\nabla T$ pour Cattaneo. Le bilan
d'énergie est $\partial_t u + \nabla\!\cdot\mathbf q = s$, avec $u = \rho c_v T$. La limite
$\tau \to 0$ restitue Fourier.

\subsection{Transformation de la source --- dérivation}

L'article donne (éq.~3) la relation entre les deux sources sans la dériver :
\[
s_{cv}(\mathbf r, t) = s_f(\mathbf r, t) + \tau\,\frac{\partial s_f(\mathbf r, t)}{\partial t}
\]
Cette relation n'est pas un choix de modélisation, c'est une conséquence algébrique. En
appliquant l'opérateur $(1 + \tau\partial_t)$ au bilan d'énergie :
\[
(1+\tau\partial_t)\big[\rho c_v\,\partial_t T\big]
+ \nabla\!\cdot\big[(1+\tau\partial_t)\mathbf q\big] = (1+\tau\partial_t)\,s
\]
Le second terme se réduit par la loi de Cattaneo, $(1+\tau\partial_t)\mathbf q = -k\nabla T$,
donc $\nabla\!\cdot[(1+\tau\partial_t)\mathbf q] = -k\nabla^2 T$. Il reste
\[
\rho c_v\,\partial_t T + \tau \rho c_v\,\partial_t^2 T - k\nabla^2 T = s + \tau\,\partial_t s
\]
soit, après division par $k$ et en utilisant $\tau/\alpha = 1/v^2$, exactement l'équation~(2)
avec $s_{cv} = s + \tau\partial_t s$. La grandeur physiquement imposée est le flux absorbé~; le
facteur $(1+\tau\partial_t)$ est le prix de l'élimination de $\mathbf q$ au profit de $T$. C'est
ce facteur qui apparaît partout ensuite sous la forme $(1+i\omega\tau)$.

\subsection{Le paramètre spectral effectif}

En transformée de Fourier temporelle, l'article écrit (éq.~5)
\[
\nabla^2 T(\mathbf r,\omega) - \sigma^2_{f/cv}\,T(\mathbf r,\omega)
= -\frac{1}{k}\,s_{f/cv}(\mathbf r,\omega)
\]
\[
\sigma_f = \left(\frac{i\omega}{\alpha}\right)^{1/2}, \qquad
\sigma_{cv} = \sqrt{\frac{i\omega}{\alpha} - \left(\frac{\omega}{v}\right)^{2}}, \qquad
s_{cv}(\mathbf r,\omega) = s_f(\mathbf r,\omega)\,[1+i\omega\tau]
\]
Le point à expliciter, que l'article ne fait pas, est que ces deux nombres d'onde sont
proportionnels :
\[
\sigma_{cv}^2 = \frac{i\omega}{\alpha} - \frac{\omega^2\tau}{\alpha}
= \frac{i\omega}{\alpha}\,(1 + i\omega\tau) = \sigma_f^2\,(1+i\omega\tau)
\]
En posant $p = i\omega$, cela s'écrit
\[
\sigma_{cv}^2 = \frac{p\,(1+\tau p)}{\alpha}
\]
C'est la substitution $p \mapsto p(1+\tau p)$ utilisée dans le modèle direct du présent travail.
L'article y parvient indépendamment, par une voie différente --- fonction de Green
tridimensionnelle axisymétrique et transformée de Hankel --- et sans en tirer la conséquence
structurelle : puisque la substitution n'affecte que $p$, la loi de Cattaneo laisse invariante
toute structure du modèle direct qui ne dépend de $p$ que par $\sigma$. C'est le résultat
négatif déjà consigné dans \code{tests/}.

Le rapport des deux paramètres spectraux,
\[
\frac{\sigma_{cv}^2}{\sigma_f^2} = 1 + i\omega\tau
\]
ne dépend que du produit $\omega\tau$. Son argument vaut $\arctan(\omega\tau)$ et sa demi-valeur
est exactement ce que la mesure de phase restitue (§~4.1).

\subsection{Dépendance en température}

L'article rappelle que $k$ et $\tau$ dépendent de la température, selon des lois de puissance
inverses $k \sim T^{-m}$ et $\tau \sim T^{-n}$. Pour le silicium, $m = 1{,}4$ et $n = 1{,}3$
(réf.~47, Glassbrenner \& Slack 1964). La limite de comportement Fourier est située à 500~K pour
Si et vers 600~K pour SiC, avec des exposants voisins. À basse température le libre parcours
moyen s'allonge et le régime devient balistique plutôt que diffusif (réf.~45, Auriault 2016).
L'hypothèse retenue est que dans un montage FDTR à température ambiante, les excursions de
température sont assez faibles pour que $k$ et $\tau$ soient constants.

Le signe des exposants est porté par la mention «~inverse power law~» et non par l'exposant
imprimé, qui est positif dans le texte. Cette ambiguïté typographique est sans conséquence sur
les calculs de l'article, puisque aucune variation en température n'y est effectuée.

% ================================================================ 3
\section{Le modèle FDTR}

\subsection{Hypothèses}

\begin{itemize}[leftmargin=1.2em,itemsep=2pt,topsep=2pt]
\item Absorption purement superficielle du pompe et de la sonde : aucune pénétration optique.
      La source est portée par un $\delta(z)$ (éq.~6).
\item Faisceaux gaussiens, pompe et sonde de même dimension pour le calcul de la phase moyennée
      (§~III~B de l'article).
\item Symétrie azimutale, problème réduit de trois à une dimension par transformée de Hankel.
\item Continuité de la température et du flux aux interfaces, par application récursive des
      conditions aux limites dans la matrice de transfert.
\item Rugosité négligée devant la longueur d'onde de la sonde (limite reconnue en conclusion).
\end{itemize}

La source spatiale est écrite (éq.~7)
\[
F(r,\omega) = \frac{F_0}{2}\,e^{-r^2/W^2}\big(1 + e^{i\omega t}\big)
\]
Cette écriture mélange les domaines temporel et fréquentiel : une fonction notée $F(r,\omega)$
ne peut contenir $e^{i\omega t}$. Le contenu utile est un profil gaussien d'échelle $W$, une
composante continue et une composante modulée à $\omega$. Par ailleurs $W$ est appelé
«~full-width half max~» alors que le profil $e^{-r^2/W^2}$ a pour largeur à mi-hauteur
$2W\sqrt{\ln 2} \simeq 1{,}665\,W$~; $W$ est le rayon à $1/e$ de l'intensité. L'écart est un
facteur $1{,}67$ sur la dimension du spot, donc un facteur $2{,}8$ sur toute grandeur
quadratique en $W$. Le cas stratifié annonce un spot de 10~\textmu m de largeur à mi-hauteur,
sans préciser laquelle des deux conventions est employée dans le calcul.

\subsection{Demi-espace}

Le problème de Green en coordonnées cylindriques (éq.~8) et la transformée de Hankel (éq.~9)
\[
g(\lambda, z; r_0, z_0; \omega) = \int_0^\infty G(r,z;r_0,z_0;\omega)\,J_0(\lambda r)\,r\,\mathrm{d}r
\]
conduisent à la solution en espace de Hankel (éq.~10)
\[
T_H^{\text{halfspace}}(\lambda, z;\omega)
= \frac{F_0 W^2}{k}\,\frac{e^{-\varepsilon(\lambda,\omega)z}}{\varepsilon(\lambda,\omega)}\,
e^{-\lambda^2 W^2/4}\,(1+i\omega\tau),
\qquad \varepsilon(\lambda,\omega) = \sqrt{\lambda^2 + \sigma^2(\omega)}
\]
puis, par transformée inverse (éq.~11),
\[
T(r;\omega) = \frac{F_0 W^2}{k}\,(1+i\omega\tau)\int_0^\infty
\frac{1}{\varepsilon(\lambda,\omega)}\,
e^{-\varepsilon(\lambda,\omega)z - \lambda^2W^2/4}\,J_0(\lambda r)\,\lambda\,\mathrm{d}\lambda
\]
On écrit $T(r, z=0,\omega) = |T|\,e^{i\phi}$. Le facteur $e^{-\lambda^2 W^2/4}$ est la
transformée de Hankel du profil gaussien~; la dépendance en $\tau$ entre par deux voies
seulement, le préfacteur $(1+i\omega\tau)$ issu de la transformation de source, et $\sigma^2$
dans $\varepsilon$. Les modèles CV et Fourier sont résolus par la même méthode et ne diffèrent
que par l'expression de $\sigma$.

L'étape intermédiaire de l'éq.~(10) telle qu'imprimée porte des facteurs
$e^{-\varepsilon}/\varepsilon$ hors de l'intégrale qui ne se retrouvent pas dans le résultat. Le
résultat final est la forme standard et c'est elle qui est employée par la suite.

\subsection{Système stratifié}

Matrice de transfert en espace de Hankel (éq.~12) :
\[
T_H^{\text{layered}}(\lambda,z;\omega) = (1+i\omega\tau)\,\frac{F_0}{k\,\varepsilon_1(\omega)}\,
W^2 e^{-\lambda^2W^2/4}\,A_0\,
\big(e^{\varepsilon_1^{cv}(\omega)z} + e^{-\varepsilon_1^{cv}(\omega)z}\big)
\]
avec $\varepsilon_i(\lambda,\omega) = \sqrt{\lambda^2+\sigma_i^2(\omega)}$, l'indice $i$
désignant la couche, et (éq.~13)
\[
A_0(\lambda,z,\omega) = \frac{e^{-\varepsilon_1 z_0}}{2k_1\varepsilon_1}
\left(\frac{e^{\varepsilon_1 z_0}[\mathcal T]_{12} + e^{-\varepsilon_1 z_0}[\mathcal T]_{22}}
{e^{\varepsilon_1 z_0}[\mathcal T]_{12} - e^{-\varepsilon_1 z_0}[\mathcal T]_{22}}\right)
\]
où $[\mathcal T]_{ij}$ est l'élément $(i,j)$ de la matrice de transfert. $A_0$ est évalué en
$z_0 = 0$, où l'expression se réduit à
$A_0 = (2k_1\varepsilon_1)^{-1}\big([\mathcal T]_{12}+[\mathcal T]_{22}\big)
\big/\big([\mathcal T]_{12}-[\mathcal T]_{22}\big)$.
Les auteurs notent que la forme de $A_0$ est voisine de l'inverse d'un coefficient de réflexion.
La température en surface du transducteur est (éq.~14)
\[
T(r,0;\omega) = 2(1+i\omega\tau)\,F_0\,k^{-1}W^2 \int_0^\infty
A_0(\lambda,\omega)\,\varepsilon_1^{-1}(\lambda,\omega)\,e^{-\lambda^2W^2/4}\,
J_0(\lambda r)\,\lambda\,\mathrm{d}\lambda
\]
La remarque structurelle importante est donnée par les auteurs : les solutions stratifiée et
demi-espace sont identiques à ceci près que la seconde comporte le coefficient de transmission
effectif $A_0$. Tout l'empilement --- conductivités, capacités, épaisseurs --- est contenu dans
$A_0$ et dans $A_0$ seul.

La phase mesurée est modélisée par une moyenne pondérée sur le spot de la sonde (réf.~51,
Cahill 2004). Avec pompe et sonde de même largeur, la grandeur pertinente est l'argument de la
moyenne complexe, soit
\[
\phi_w(\omega) = \arg\!\left[(1+i\omega\tau)\int_0^\infty
\frac{A_0(\lambda,\omega)}{\varepsilon_1(\lambda,\omega)}\,e^{-\lambda^2W^2/2}\,
\lambda\,\mathrm{d}\lambda\right]
\]
le second facteur gaussien provenant du recouvrement pompe--sonde. L'article écrit «~weighted
average of the phase~», c'est-à-dire la moyenne de la phase, qui n'est pas l'argument de la
moyenne. Les deux grandeurs diffèrent dès que la phase varie sur le spot, ce qui est
précisément le cas dans les figures 4(c) et 4(d), où $\phi$ passe de $0$ à $-80^\circ$ sur
$r/W \in [0,2]$. L'écart n'est pas quantifié dans l'article.

\subsection{Ce que le modèle ne contient pas}

Recensement, par lecture exhaustive des équations (1) à (15) et des paramètres annoncés.

\smallskip
{\small
\begin{tabularx}{\linewidth}{@{}L{0.27\linewidth} L{0.22\linewidth} X@{}}
\toprule
\textbf{Élément} & \textbf{Présence} & \textbf{Conséquence} \\
\midrule
Conductance thermique d'interface $G$ & absente --- continuité de $T$ et $q$ &
le compétiteur principal de $\tau$ dans un ajustement Au/SiC est retiré de l'analyse \\
Pertes en face avant & absentes & sans effet aux fréquences considérées \\
Capacité localisée du transducteur & non --- l'or est une couche à $k,\alpha$ propres &
cohérent, plus détaillé que le modèle localisé usuel \\
Pénétration optique & négligée & hypothèse standard pour un transducteur métallique \\
Anisotropie & absente & signalée en conclusion comme extension possible \\
Modèle de bruit & absent & $S_\tau$ n'est jamais converti en incertitude \\
Matrice de covariance, matrice de Fisher & absentes & l'identifiabilité n'est pas abordée \\
Ajustement, inversion, données synthétiques & absents &
aucune estimation de $\tau$ n'est effectuée dans l'article \\
Terme non local, Guyer--Krumhansl & absents &
le modèle reste strictement Cattaneo~; DPL mentionné, non utilisé \\
\bottomrule
\end{tabularx}}

\smallskip
L'absence de conductance d'interface est la plus lourde des neuf. Dans une expérience FDTR
réelle sur Au/SiC, $G$ est un paramètre d'ajustement dominant, et il est connu pour être
fortement corrélé à la conductivité du substrat. Son absence retire du calcul de sensibilité le
paramètre dont le vecteur de sensibilité est le plus susceptible d'être colinéaire à celui de
$\tau$.

% ================================================================ 4
\section{Résultat analytique central --- la phase du demi-espace}

\subsection{Forme fermée}

L'article donne la limite haute fréquence pour le système stratifié (éq.~15) :
\[
\phi \simeq -\frac{\pi}{4} + \frac{1}{2}\arctan(\omega\tau)
+ \arctan\!\left(\frac{\operatorname{Im}A_0}{\operatorname{Re}A_0}\right)
\]
et précise que le troisième terme n'existe pas pour le demi-espace. La forme du demi-espace se
retrouve directement. En régime unidimensionnel, l'intégrale de l'éq.~(11) est dominée par
$\lambda \ll |\sigma|$, donc $\varepsilon \simeq \sigma$ et
\[
T \propto \frac{1+i\omega\tau}{\sigma_{cv}}
= \frac{1+i\omega\tau}{\sqrt{i\omega/\alpha}\,\sqrt{1+i\omega\tau}}
= \sqrt{1+i\omega\tau}\;\sqrt{\frac{\alpha}{i\omega}}
\]
d'où
\begin{encadre}
\[
\phi(\omega) = -\frac{\pi}{4} + \frac{1}{2}\arctan(\omega\tau)
\]
\end{encadre}

Trois lectures de ce résultat.

\textbf{Première.} La dépendance en $\tau$ de la phase mesurée est \textit{la moitié de
l'argument du rapport des paramètres spectraux} établi au §~2.3. La grandeur $1+i\omega\tau$ est
le seul véhicule de $\tau$ dans l'observable, et la mesure n'en voit que la demi-phase.

\textbf{Deuxième.} Les deux asymptotes sont $-\pi/4$ pour $\omega\tau \ll 1$ et $0$ pour
$\omega\tau \gg 1$. La première est la phase de Fourier en régime unidimensionnel~; la seconde
correspond au régime d'onde non amortie, où $\sigma_{cv}^2 \simeq -\omega^2/v^2$ donne
$\sigma_{cv} \simeq i\omega/v$, nombre d'onde purement imaginaire. Le déphasage de $-90^\circ$
apporté par $1/\sigma_{cv}$ est alors exactement compensé par le $+90^\circ$ de
$(1+i\omega\tau)$. L'annulation est celle que les auteurs décrivent : «~in the high-frequency
limit, the $\arctan(\omega\tau)$ asymptotically approaches $\pi/2$, and this leads to
cancellation of the Fourier phase~».

\textbf{Troisième.} La phase CV n'est pas monotone en fréquence pour un spot fini. À basse
fréquence la longueur de diffusion excède le spot, la réponse est quasi stationnaire
tridimensionnelle et $\phi \to 0$. Le régime unidimensionnel s'installe ensuite et tire $\phi$
vers $-45^\circ$. Puis le terme en $\arctan(\omega\tau)$ la ramène vers $0$. La figure 3(a)
montre ce minimum intermédiaire, situé autour de $\omega\tau \simeq 0{,}2$, à une valeur
d'environ $-32^\circ$. La courbe de Fourier, elle, est monotone décroissante et sature à
$-45^\circ$.

La forme fermée unidimensionnelle donnerait, en $\omega\tau = 0{,}2$, une phase de
$-39{,}3^\circ$, contre environ $-32^\circ$ lus sur la figure. L'écart de $7^\circ$ mesure la
contribution tridimensionnelle résiduelle du spot à cette fréquence : le minimum de la courbe
n'est donc pas atteint dans le domaine de validité de la forme fermée. Ce constat est à
rapprocher du §~5.4, où la localisation de l'extremum de $S_\tau$ est en revanche conforme à la
prédiction unidimensionnelle.

\subsection{Le seuil $\omega\tau = 1$}

La phase du demi-espace atteint la valeur médiane entre ses deux asymptotes,
$-\pi/8 = -22{,}5^\circ$, exactement en $\omega\tau = 1$, puisque $\arctan(1) = \pi/4$. Le seuil
n'est donc pas qualitatif : c'est le centre exact de la transition de phase, et c'est le même
seuil que celui du critère spectral du présent travail.

La dérivée logarithmique de la phase par rapport à $\tau$ se calcule sur la forme fermée :
\[
\frac{\mathrm{d}\phi}{\mathrm{d}\ln\tau} = \tau\,\frac{\partial\phi}{\partial\tau}
= \frac{1}{2}\,\frac{\omega\tau}{1+(\omega\tau)^2}
\]
Cette fonction est maximale en $\omega\tau = 1$, où elle vaut $1/4$ radian, soit
$14{,}32^\circ$ par facteur $e$ sur $\tau$, ou $32{,}98^\circ$ par décade. Elle décroît comme
$\omega\tau/2$ en deçà du seuil et comme $1/(2\omega\tau)$ au-delà.

\smallskip
{\small\centering
\begin{tabular}{@{}rrr@{}}
\toprule
$\omega\tau$ & $\mathrm{d}\phi/\mathrm{d}\ln\tau$ (deg) & par décade (deg) \\
\midrule
0,1 & 2,84 & 6,53 \\
0,3 & 7,88 & 18,16 \\
1 & 14,32 & 32,98 \\
3 & 8,59 & 19,79 \\
10 & 2,84 & 6,53 \\
100 & 0,286 & 0,660 \\
\bottomrule
\end{tabular}\par}

\smallskip
La décroissance en $1/(2\omega\tau)$ au-delà du seuil est un point que l'article n'énonce pas et
qui contredit l'intuition selon laquelle monter en fréquence améliore indéfiniment la mesure de
$\tau$ : la sensibilité est maximale au seuil et se dégrade de part et d'autre, symétriquement
en échelle logarithmique.

\subsection{Le cas stratifié et l'enchevêtrement par $A_0$}

Pour le système stratifié, le troisième terme de l'éq.~(15) réintroduit tout l'empilement. Les
auteurs le formulent en ces termes : le troisième terme n'existe pas pour le demi-espace, «~so
the value of the relaxation time could be determined through the phase alone~». La proposition
contraposée, non écrite, est que pour le système stratifié la phase seule ne suffit pas ---
puisque $\tau$ et $A_0$ y entrent additivement dans la même observable et que $A_0$ contient les
conductivités, les capacités et les épaisseurs par la matrice de transfert.

Trois régimes sont identifiés dans l'espace des fréquences :

\begin{itemize}[leftmargin=1.2em,itemsep=2pt,topsep=2pt]
\item $\omega \ll \tau^{-1}$ : la phase est approximativement constante, sans différence
      significative entre modèles~; la dynamique est gouvernée par Fourier.
\item $\omega \sim \tau^{-1}$ : régime de transition, où $\tau$ devient pertinent~; pour le
      système stratifié, ce sont les parties réelle et imaginaire de $A_0$, qui contiennent le
      nombre d'onde thermique $\sigma$, qui fixent le comportement.
\item $\omega \gg \tau^{-1}$ : éq.~(15).
\end{itemize}

La dépendance de $\phi$ en l'épaisseur de la couche entre par la matrice de transfert, ce que
les auteurs qualifient de «~not straightforward~», en constatant empiriquement (§~6.2) qu'une
épaisseur de SiC de 1~\textmu m maximise l'écart CV--Fourier sans relation directe apparente
avec l'épaisseur. Une explication quantitative s'obtient par confrontation d'échelles. La
longueur de diffusion à la fréquence de seuil vaut
\[
\mu\big(\omega = \tau^{-1}\big) = \sqrt{\frac{2\alpha}{\omega}} = \sqrt{2\alpha\tau}
= \sqrt{2}\;v\tau
\]
c'est-à-dire, à un facteur $\sqrt 2$ près, la distance parcourue par l'onde thermique pendant un
temps de relaxation. La signature CV est maximale lorsque la couche est résolue exactement à la
transition, soit $d \simeq \sqrt{2\alpha\tau}$. Avec des valeurs de littérature pour le SiC
($k \simeq 370$~W\,m$^{-1}$K$^{-1}$, $\rho c_v \simeq 2{,}2\times10^6$~J\,m$^{-3}$K$^{-1}$, donc
$\alpha \simeq 1{,}7\times10^{-4}$~m$^2$\,s$^{-1}$) et $\tau = 10^{-8}$~s, il vient
$v \simeq 130$~m\,s$^{-1}$, $v\tau \simeq 1{,}3$~\textmu m et
$\sqrt{2\alpha\tau} \simeq 1{,}8$~\textmu m. Pour une couche de 10~\textmu m, la couche paraît
semi-infinie au seuil et le système se comporte comme le demi-espace~; pour 0,1~\textmu m, la
couche est thermiquement mince au seuil et le substrat de Si, traité comme matériau de Fourier,
domine. L'optimum observé à 1~\textmu m est donc la coïncidence de l'épaisseur avec la longueur
balistique $v\tau$.

Une troisième échelle achève de fixer la géométrie du problème au seuil. En régime d'onde,
$\sigma_{cv} \simeq i\omega/v$, et la longueur d'onde thermique vaut $\Lambda = 2\pi v/\omega$,
soit $2\pi v\tau \simeq 8{,}1$~\textmu m en $\omega = \tau^{-1}$. Trois longueurs coexistent
donc à la transition : $\Lambda \simeq 8$~\textmu m, la largeur du spot annoncée à
10~\textmu m, et l'épaisseur de la couche à 1~\textmu m. La première et la deuxième étant du
même ordre, l'hypothèse unidimensionnelle n'est pas acquise au seuil pour le système stratifié,
ce qui justifie que le comportement y soit gouverné par $A_0$ et non par la forme fermée de
l'éq.~(15).

Les valeurs de $\alpha$ employées dans ce paragraphe sont de littérature, la table de l'article
figurant dans le matériel supplémentaire non consulté.

% ================================================================ 5
\section{La définition de la sensibilité}

\subsection{Définition employée}

L'article écrit, à la fin de la section III~A :
\[
S_\tau = \frac{\mathrm{d}\phi}{\mathrm{d}\ln\tau}
\]
en renvoyant à la référence 51 (Cahill, \textit{Rev. Sci. Instrum.} \textbf{75}, 5119, 2004).
Puis : «~We observe that FDTR shows significant sensitivity for frequencies at the $\tau^{-1}$
threshold~; therefore, it will be possible to fit for the time-relaxation parameter in this
frequency range.~»

C'est la seule grandeur de type métrologique calculée dans l'article, et elle est calculée pour
les trois configurations : figure 3(b) pour le demi-espace, figure 4(e) pour Au/SiC/Si,
figure 6(e) pour le miroir de Bragg.

Deux remarques sur la définition elle-même. D'une part, le coefficient de sensibilité de
Cahill (2004) est construit pour le signal de rapport $-V_{in}/V_{out}$ en TDTR~; sa
transposition à la phase en FDTR est légitime mais n'est pas discutée. D'autre part, $S_\tau$
est une dérivée logarithmique en $\tau$ mais pas en $\phi$ : elle est donc homogène à un angle,
et sa valeur numérique dépend de l'unité d'angle et de la base du logarithme retenues. Ni l'une
ni l'autre n'est spécifiée sur les axes des figures 3(b), 4(e) et 6(e), qui portent
«~$S_\tau$~» sans unité.

\subsection{Rappel --- sensibilité et identifiabilité}

Soit $\boldsymbol\theta$ le vecteur des paramètres inconnus,
$\phi(\omega_m;\boldsymbol\theta)$ le modèle direct, $\mathbf J$ la matrice jacobienne
$J_{mi} = \partial\phi(\omega_m)/\partial\theta_i$ et $\boldsymbol\Sigma$ la covariance du bruit
de phase. La borne de Cramér--Rao sur l'écart-type d'estimation du paramètre $i$ est
\[
\sigma^2_{\theta_i}
= \Big[\big(\mathbf J^{\!\top}\boldsymbol\Sigma^{-1}\mathbf J\big)^{-1}\Big]_{ii}
\]
Dans le cas d'un paramètre unique et d'un bruit blanc d'écart-type $\sigma_\phi$, cela se réduit
à
\[
\sigma^2_{\ln\tau} = \frac{\sigma_\phi^2}{\sum_m S_\tau(\omega_m)^2}
\]
et une grande valeur de $|S_\tau|$ borne effectivement l'incertitude. Dès que
$\dim\boldsymbol\theta > 1$, l'inversion de
$\mathbf J^{\!\top}\boldsymbol\Sigma^{-1}\mathbf J$ fait intervenir les corrélations, et
l'élément diagonal $ii$ n'est plus contrôlé par la norme de la colonne $\mathbf J_{:,i}$ mais
par sa composante orthogonale au sous-espace engendré par les autres colonnes. En notant
$\mathbf J_{:,\tau}^{\perp}$ cette composante,
\[
\sigma_{\ln\tau} = \frac{\sigma_\phi}{\big\|\mathbf J_{:,\tau}^{\perp}\big\|}
\]
de sorte que $\sigma_{\ln\tau}$ peut être arbitrairement grande avec
$\|\mathbf J_{:,\tau}\|$ arbitrairement grande, si l'angle entre $\mathbf J_{:,\tau}$ et le
sous-espace des autres colonnes tend vers zéro. La grandeur déterminante est un angle, non une
norme.

C'est exactement la distinction que Krapez \& Rigollet opposent à Guo \textit{et al.} --- voir
\code{notes/\allowbreak krapez\_2017.md}, §~6 : une dérivée du signal grande n'établit pas qu'un paramètre
soit identifiable, seule la structure de corrélation le fait.

\subsection{L'inférence de l'article}

L'énoncé de l'article est de la forme «~$|S_\tau|$ est significative au seuil, \textbf{donc} il
sera possible d'ajuster $\tau$ dans cette bande~». Le \textit{donc} n'est pas établi, et ne peut
pas l'être à partir de $S_\tau$ seul, pour les raisons du §~5.2. Aucun des éléments suivants
n'apparaît dans l'article : jacobienne complète, matrice de corrélation, borne de
Cramér--Rao, écart-type de phase injecté, ajustement sur données synthétiques.

Les paramètres maintenus fixes pendant le calcul de $S_\tau$, pour le cas Au/SiC/Si, sont au
minimum : $k$ et $\rho c_v$ de l'or, du SiC et du Si (six valeurs, table supplémentaire),
l'épaisseur du transducteur (100~nm), l'épaisseur de SiC (1~\textmu m), la dimension du spot
(10~\textmu m), la puissance absorbée (1~mW). Aucun n'est varié conjointement avec $\tau$. À
quoi s'ajoute, dans une expérience réelle, la conductance d'interface Au/SiC, absente du modèle
(§~3.4), et la phase de référence du montage.

Il faut noter que la localisation de $|S_\tau|$ maximal aggrave le problème plutôt qu'elle ne le
résout. Dans le cas stratifié, l'extremum de la figure 4(e), $|S_\tau| \simeq 58$ près de
$10^9$ sur l'abscisse, coïncide avec l'excursion brutale de phase de la figure 4(f) --- chute à
environ $-70^\circ$ puis remontée à $+5^\circ$. Cette excursion est une quasi-résonance de
$A_0$. En un tel point, toutes les sensibilités croissent simultanément, puisqu'elles sont
toutes dominées par la dérivée du même dénominateur de $A_0$. C'est le régime où la colinéarité
est maximale, et donc le plus défavorable du point de vue de l'identifiabilité, alors qu'il est
le plus favorable du point de vue de la norme de la sensibilité.

\subsection{Contrôle numérique de $S_\tau$}

La forme fermée du §~4.2 fournit un contrôle indépendant de la figure 3(b).

\textbf{Position.} L'extremum analytique est en $\omega\tau = 1$, soit $\omega = 10^{8}$ pour
$\tau = 10^{-8}$~s. L'extremum de la figure 3(b) est situé aux environs de $10^{8}$ sur
l'abscisse. La position concorde, et elle fixe la lecture de l'abscisse (§~7.1).

\textbf{Amplitude.} Le maximum analytique est $14{,}32^\circ$ par facteur $e$, soit
$32{,}98^\circ$ par décade. La figure 3(b) donne $|S_\tau| \simeq 30$. Le rapport entre
l'amplitude tracée et la borne analytique par facteur $e$ est de $2{,}1$, proche de
$\ln 10 = 2{,}303$. Deux lectures sont possibles, sans que le texte permette de trancher : soit
la grandeur tracée est $\mathrm{d}\phi/\mathrm{d}\log_{10}\tau$ et non
$\mathrm{d}\phi/\mathrm{d}\ln\tau$, soit le maximum tracé sort du domaine de validité de la
forme asymptotique unidimensionnelle. Le facteur $2{,}303$ entre les deux conventions se
répercute directement sur toute conversion de $S_\tau$ en incertitude.

\textbf{Signe.} La dérivée analytique $\tfrac12\,\omega\tau/(1+(\omega\tau)^2)$ est strictement
positive : augmenter $\tau$ relève la phase vers $0$. La figure 3(b) montre $S_\tau$ négatif,
d'extremum $\simeq -30$. Sous la convention de phase de la figure 3(a), où la courbe de Fourier
sature à $-45^\circ$, la grandeur tracée est donc $-\mathrm{d}\phi/\mathrm{d}\ln\tau$, ou la
phase employée dans la définition est le retard $-\phi$. L'article ne le précise pas. Pour le
cas stratifié, la figure 4(e) présente les deux signes, ce qui est attendu : le terme en $A_0$
de l'éq.~(15) a une dérivée en $\tau$ de signe variable.

\textbf{Amplitudes comparées des trois configurations.} $|S_\tau|_{\max} \simeq 30$ pour le
demi-espace, $\simeq 58$ pour Au/SiC/Si, $\simeq 0{,}33$ pour le miroir de Bragg. L'étendue
couvre plus de deux ordres de grandeur. La «~haute sensibilité~» annoncée dans le résumé n'est
donc pas une propriété de la technique FDTR, mais de la configuration et de la portion de l'axe
$\omega\tau$ effectivement atteinte.

\textbf{Ce qu'une conversion naïve donnerait.} L'article invoque une résolution de phase de
quelques millidegrés (réf.~40) sans jamais la combiner à $S_\tau$. La combinaison, pour un
unique point de mesure au seuil et $\tau$ seul inconnu, donnerait
$\delta\ln\tau \simeq \delta\phi/|S_\tau| \simeq 0{,}003/30 = 10^{-4}$, soit une incertitude
relative de $10^{-2}$~\% sur un paramètre dont l'article rappelle par ailleurs qu'aucun modèle
théorique ni procédure expérimentale n'en fixe la valeur. L'invraisemblance du résultat est le
diagnostic : c'est le calcul à un paramètre qui n'est pas la grandeur pertinente.

% ================================================================ 6
\section{Résultats numériques et configurations}

\subsection{Demi-espace SiC}

$\tau_{SiC} = 10^{-8}$~s, valeur déclarée «~of the order but not equal to reported phonon
relaxation times~», en renvoyant à Ezzahri \& Joulain (réf.~50), article théorique sur la
conductivité thermique dynamique des cristaux semiconducteurs. Il s'agit donc d'un choix de
plausibilité, non d'une mesure. L'article rappelle en introduction qu'aucun modèle théorique ni
procédure expérimentale n'existe pour établir $\tau$ pour un matériau donné, et ne distingue pas
le temps de relaxation de Cattaneo de la durée de vie des phonons autrement que par cette
réserve.

Figure 2 : cartes de $|\Delta T_{CV}|$ et $|\Delta T_F|$ en fonction de l'abscisse fréquentielle
($10^5$ à $10^9$) et de $r/W_0 \in [0,5]$, échelle jusqu'à $0{,}3$~K~; coupe spectrale (c) et
coupe spatiale (d). Sur la coupe spectrale, $|\Delta T|$ passe de $0{,}175$~K à $10^5$ à environ
$0{,}02$~K (CV) contre $0{,}005$~K (Fourier) à $10^9$. Sur la coupe spatiale, $0{,}07$~K (CV)
contre $0{,}056$~K (Fourier) au centre du spot. Les deux modèles ne diffèrent qu'au voisinage
immédiat du maximum d'intensité laser~; au-delà, l'écart devient négligeable.

Ce résultat sur le module est cohérent avec la structure du §~2.3 : le module de la réponse est
peu discriminant, l'information sur $\tau$ étant portée par la phase. C'est la même conclusion
que celle du présent travail sur l'amplitude contre la phase.

Figure 3 : phase pondérée (a) et $S_\tau$ (b), $\tau = 10^{-8}$~s. Analysée aux §~4.1 et §~5.4.

\subsection{Au/SiC/Si}

Configuration : couche de SiC de 1~\textmu m sur substrat de Si, transducteur d'or de 100~nm,
spot gaussien de 10~\textmu m de largeur à mi-hauteur, puissance transférée au système 1~mW,
$\tau_{SiC} = 10^{-8}$~s. L'or et le silicium sont traités comme matériaux de Fourier, de sorte
que seul le SiC possède un temps de réponse non négligeable.

Figure 4 : cartes $|\Delta T_{CV}|$, $|\Delta T_F|$ (jusqu'à $0{,}3$~K), $\phi_{CV}$, $\phi_F$
(de $0$ à $-80^\circ$) sur $r/W \in [0,2]$~; $S_\tau(\omega)$ en (e), lobe positif
$\simeq +28$ vers $10^8$ et excursion négative $\simeq -58$ vers $10^9$~; phase pondérée en (f).
Le module est très voisin entre les deux modèles~; c'est près de $\tau^{-1}$ qu'apparaît une
forte région de retard de phase absente du modèle de Fourier.

Figure 5(a) : phase pondérée, $\tau \in \{1, 3, 6, 10\}\times 10^{-9}$~s à
$d_{SiC} = 10^{-6}$~m fixée. Énoncés associés : dans cette bande, les temps de relaxation
inférieurs à $10^{-9}$~s sont pratiquement identiques et représentent la limite de Fourier du
système~; la phase CV tend vers celle de Fourier quand $\tau$ tend vers zéro~; des écarts
restent apparents jusqu'à $3\times10^{-9}$~s, mais pour cette valeur les différences se situent
hors de la gamme de fréquences des montages FDTR courants.

Figure 5(b) : $d_{SiC} \in \{10^{-5}, 10^{-6}, 10^{-7}\}$~m à $\tau = 10^{-8}$~s fixé, traits
pleins pour CV, pointillés pour Fourier, abscisse de $10^4$ à $10^9$. Dépendance non triviale en
épaisseur~; les écarts apparaissent à plus haute fréquence pour la couche la plus mince~; aucune
relation directe apparente entre épaisseur et écart, $d_{SiC} = 10^{-6}$~m produisant les écarts
les plus grands et constituant le cas idéal. Explication quantitative au §~4.3.

\subsection{Miroir de Bragg thermique}

Configuration : transducteur d'or de 100~nm, cellule unité formée d'une couche de derme de
1~\textmu m et d'une couche d'épiderme de 1~\textmu m, de constantes de temps respectives 1~s et
2~s, cellule répétée trois fois, substrat de Si. Paramètres repris de la réf.~11 (Camacho de la
Rosa \textit{et al.}, \textit{Energies} \textbf{14}, 7452, 2021) pour permettre la comparaison.
Spectre d'excitation de 0 à 40~Hz.

Figure 6 : cartes de température (jusqu'à environ 7~mK) et de phase~; $S_\tau(\omega)$ en (e),
entre $+0{,}15$ et $-0{,}33$~; phase pondérée en (f), atteignant environ $+85^\circ$ pour CV
contre $+6^\circ$ pour Fourier à 40~Hz. Observations : moins d'échauffement de la surface du
transducteur pour CV à basse fréquence~; formation d'un gap autour de 30~Hz, attribué aux
interférences d'ondes thermiques entre couches~; à 30~Hz, creux de la distribution de
température même pour $r/W < 1$, cohérent avec une forte transmissivité thermique dans cette
fenêtre~; à 40~Hz, température comparable pour $r/W > 1$ et $r/W < 1$, indiquant que l'énergie
thermique n'est pas aisément transmise selon $z$.

Deux réserves de lecture. D'une part, la phase est ici \textbf{positive}, jusqu'à $+85^\circ$,
alors qu'elle est négative dans toutes les autres figures. Pour une réponse causale avec la
convention $T \propto e^{i\omega t}$, la température de surface est en retard sur le flux et la
phase doit être négative. Les panneaux 6(c), 6(d) et 6(f) portent donc soit le retard $-\phi$,
soit une phase déroulée à travers le gap. L'article ne signale pas le changement.

D'autre part, la phrase introductive de la section III~C est contradictoire telle qu'imprimée :
«~Although the frequency range of typical FDTR systems is in the order of tens of hertz and is
outside the range of typical working conditions of FDTR~». La lecture cohérente est que c'est la
gamme de fréquences \textit{de ce système}, fixée par $1/\tau$ avec $\tau$ de l'ordre de la
seconde, qui est de quelques dizaines de hertz, et qui est hors des conditions habituelles de
fonctionnement du FDTR (10~kHz -- 200~MHz). La configuration est donc explicitement présentée
comme non réalisable dans un montage FDTR standard, et retenue pour la seule comparaison avec le
travail antérieur.

% ================================================================ 7
\section{Bande de fréquences}

\subsection{Ambiguïté d'abscisse}

Toutes les figures portent en abscisse «~$\omega$ [Hz]~». L'étiquette est dimensionnellement
incohérente : $\omega$ est une pulsation, en rad\,s$^{-1}$, et $\omega = 2\pi f$ avec $f$ en Hz.
Deux lectures sont possibles, séparées par un facteur $2\pi \simeq 6{,}28$.

Trois éléments internes départagent en faveur de l'abscisse en rad\,s$^{-1}$.

\begin{enumerate}[leftmargin=1.6em,itemsep=2pt,topsep=2pt]
\item L'extremum de $S_\tau$ en figure 3(b) est situé vers $10^8$, et l'extremum analytique est
      en $\omega\tau = 1$, soit $\omega = \tau^{-1} = 10^8$~rad\,s$^{-1}$ pour
      $\tau = 10^{-8}$~s (§~5.4). Si l'abscisse portait $f$, l'extremum devrait se trouver à
      $f = 1/(2\pi\tau) = 1{,}59\times10^7$.
\item Le texte identifie systématiquement le seuil à la valeur $\tau^{-1}$ lue directement sur
      l'abscisse.
\item L'étendue des abscisses, $10^4$ à $10^9$, correspond en fréquence à
      $1{,}6$~kHz~--~$160$~MHz, soit très exactement la bande 10~kHz~--~200~MHz que l'article
      annonce lui-même en introduction comme gamme usuelle de modulation FDTR (réf.~39 et 40).
\end{enumerate}

Sous cette lecture, la fréquence de modulation au seuil est $f = 1/(2\pi\tau)$, et non
$\tau^{-1}$.

\subsection{Table des seuils}

Bande FDTR retenue par l'article : 10~kHz à 200~MHz, la limite supérieure dépendant du rapport
signal sur bruit (réf.~39, Regner \textit{et al.} 2013~; réf.~40, Kirsch \textit{et al.} 2024).

\smallskip
{\small\centering
\begin{tabular}{@{}rrrl@{}}
\toprule
$\tau$ (s) & $\tau^{-1}$ (rad\,s$^{-1}$) & $f_{\text{seuil}} = 1/(2\pi\tau)$ &
dans 10~kHz -- 200~MHz \\
\midrule
$10^{-6}$            & $10^{6}$             & 159~kHz   & oui \\
$10^{-7}$            & $10^{7}$             & 1,59~MHz  & oui \\
$10^{-8}$            & $10^{8}$             & 15,9~MHz  & oui \\
$3\times10^{-9}$     & $3{,}3\times10^{8}$  & 53,1~MHz  & oui \\
$10^{-9}$            & $10^{9}$             & 159~MHz   & limite \\
$8{,}0\times10^{-10}$& $1{,}25\times10^{9}$ & 199~MHz   & plafond \\
$10^{-10}$           & $10^{10}$            & 1,59~GHz  & non \\
$10^{-12}$           & $10^{12}$            & 159~GHz   & non \\
\bottomrule
\end{tabular}\par}

\smallskip
Le plus petit temps de relaxation dont le centre de transition soit atteignable à 200~MHz est
\begin{encadre}
\[
\tau_{\min} = \frac{1}{2\pi f_{\max}} = 8{,}0\times10^{-10}\ \text{s}
\]
\end{encadre}
Cette valeur est indépendante du matériau et ne dépend que du plafond instrumental. Elle est
cohérente avec l'énoncé de l'article selon lequel les temps inférieurs à $10^{-9}$~s sont
indistinguables de Fourier dans la bande tracée (§~6.2), ce qui constitue un quatrième argument
en faveur de la lecture du §~7.1.

\subsection{L'énoncé des auteurs}

L'article ne revendique pas la mesure de temps de relaxation arbitrairement courts. Il énonce,
pour $\tau = 3\times10^{-9}$~s : «~While the analytical analysis is valid, FDTR capabilities will
need to increase to reach this frequency range.~» Et pour la configuration de Bragg, que la
gamme requise est hors des conditions de fonctionnement habituelles du FDTR (§~6.3).

Il y a donc, dans un même article, une revendication de haute sensibilité au résumé et à la
conclusion, et la reconnaissance dans le corps du texte que la bande où cette sensibilité
s'exprime n'est pas atteinte par l'instrumentation existante pour
$\tau \lesssim 3\times10^{-9}$~s. Les deux énoncés ne sont pas contradictoires --- l'un porte
sur une dérivée, l'autre sur une plage d'accès --- mais leur coexistence interdit de citer le
premier sans le second.

% ================================================================ 8
\section{Écart entre le corps de l'article et sa conclusion}

La conclusion (§~IV de l'article) énonce quatre propositions qui ne sont pas établies par le
corps du texte.

\smallskip
{\small
\begin{tabularx}{\linewidth}{@{}X X@{}}
\toprule
\textbf{Énoncé de la conclusion} & \textbf{État dans le corps de l'article} \\
\midrule
«~demonstrates a robust methodology for determining relaxation times~» &
aucune méthodologie d'estimation n'est spécifiée : ni fonction de coût, ni algorithme, ni plan
de fréquences, ni traitement des paramètres de nuisance \\
«~we effectively captured high-resolution transient thermal responses~» &
aucune donnée, expérimentale ou synthétique, n'est produite~; toutes les figures sont des
calculs de modèle direct \\
«~enabling precise fitting to a non-Fourier heat conduction model~» &
aucun ajustement n'est effectué \\
«~allowing for accurate quantification of relaxation times~» &
aucune quantification, donc aucune incertitude \\
\bottomrule
\end{tabularx}}

\smallskip
Ce qui est effectivement établi est plus restreint et n'en est pas moins utilisable : la forme
de la réponse en phase sous CV pour trois géométries, la localisation du régime de transition en
$\omega \sim \tau^{-1}$, le calcul du coefficient $S_\tau$ dans ces trois géométries, et la
constatation que le module de la réponse est peu discriminant tandis que la phase l'est.

Une limite est reconnue explicitement : la rugosité est supposée négligeable devant la longueur
d'onde de la sonde, alors qu'en pratique elle doit être mesurée car elle peut induire un
changement de phase sans lien avec un comportement non-Fourier (réf.~52 et 53). La remarque est
plus lourde qu'elle n'en a l'air, puisqu'un décalage de phase systématique est exactement la
forme que prend la signature revendiquée. Elle s'ajoute aux deux autres sources de décalage de
phase non modélisées : la conductance d'interface absente du modèle (§~3.4) et la phase de
référence du montage.

% ================================================================ 9
\section{Points non tranchés et obscurités relevées}

Liste exhaustive des points que la lecture du seul texte publié ne permet pas de fixer.

\begin{enumerate}[leftmargin=1.7em,itemsep=3pt,topsep=2pt]
\item \textbf{Unité et base du logarithme de $S_\tau$.} Facteur $\ln 10 = 2{,}303$ d'écart entre
      $\mathrm{d}\phi/\mathrm{d}\ln\tau$ et $\mathrm{d}\phi/\mathrm{d}\log_{10}\tau$ (§~5.4).
\item \textbf{Signe de $S_\tau$.} Le signe tracé est opposé à celui de la dérivée analytique
      sous la convention de phase des figures (§~5.4).
\item \textbf{Convention d'abscisse.} $\omega$ en rad\,s$^{-1}$ ou $f$ en Hz~; facteur $2\pi$
      (§~7.1). L'évidence interne tranche en faveur de rad\,s$^{-1}$, sans que le texte le dise.
\item \textbf{Convention de phase de la figure 6.} Phase positive jusqu'à $+85^\circ$, contre
      phase négative partout ailleurs (§~6.3).
\item \textbf{Définition de $W$.} Rayon à $1/e$ ou largeur à mi-hauteur~; facteur $1{,}665$
      (§~3.1).
\item \textbf{Moyenne de la phase ou argument de la moyenne.} Les deux diffèrent là où la phase
      varie sur le spot, c'est-à-dire précisément dans la région exploitée (§~3.3).
\item \textbf{Propriétés thermiques employées.} Table 1 du matériel supplémentaire, non
      consultée~; aucun des résultats numériques n'est reproductible à partir du seul corps de
      l'article.
\item \textbf{Expression de la matrice de transfert.} Matériel supplémentaire également. Le
      préfacteur $\varepsilon_1^{-1}$ apparaît à la fois dans $A_0$ (éq.~13) et dans l'intégrande
      de l'éq.~(14)~; la cohérence des deux écritures demande la dérivation complète.
\item \textbf{Étape intermédiaire de l'éq.~(10).} Facteurs $e^{-\varepsilon}/\varepsilon$ hors
      intégrale, absents du résultat final (§~3.2).
\item \textbf{Écriture de l'éq.~(7).} Mélange des domaines temporel et fréquentiel (§~3.1).
\item \textbf{Phrase introductive de la section III~C.} Contradictoire telle qu'imprimée~;
      lecture cohérente proposée au §~6.3.
\item \textbf{Disponibilité des données.} «~available from the corresponding author upon
      reasonable request~». Aucun code, aucun jeu déposé.
\end{enumerate}

Aucun de ces points n'affecte la validité de la structure analytique des §~2 à~4, qui a été
recalculée indépendamment. Les points 1, 2 et 3 affectent toute conversion numérique de
$S_\tau$ en incertitude~; les points 7 et 8 affectent la reproductibilité des figures.

% ================================================================ 10
\section{Position par rapport au présent travail}

\subsection{Ce que l'article confirme}

\textbf{La substitution spectrale.} $\sigma_{cv}^2 = p(1+\tau p)/\alpha$ (§~2.3) est la
substitution implantée dans \code{src/\allowbreak forward\_model.py}, obtenue ici par une voie
indépendante --- Green tridimensionnel et Hankel plutôt que quadripôles unidimensionnels. La
concordance vaut vérification croisée du modèle direct.

\textbf{Le seuil.} Le seuil $\omega\tau = 1$ apparaît dans l'article comme frontière des trois
régimes, et analytiquement comme centre exact de la transition de phase et maximum exact de
$\mathrm{d}\phi/\mathrm{d}\ln\tau$ (§~4.2). C'est le même seuil que celui du critère spectral du
présent travail, atteint par une autre route.

\textbf{La phase contre le module.} L'article établit numériquement, dans les trois
configurations, que le module de la réponse discrimine mal les deux modèles tandis que la phase
les discrimine.

\subsection{Ce que l'article laisse ouvert}

\textbf{L'identifiabilité.} L'article calcule $\partial\phi/\partial\ln\tau$~; il ne calcule ni
$\big[(\mathbf J^{\!\top}\boldsymbol\Sigma^{-1}\mathbf J)^{-1}\big]_{\tau\tau}$, ni aucune
corrélation. La question de savoir si $\tau$ se sépare des autres paramètres n'est pas posée.

\textbf{L'enchevêtrement par $A_0$.} Les auteurs constatent que pour le demi-espace $\tau$ peut
être déterminé par la phase seule parce que le terme en $A_0$ est absent, et n'énoncent pas la
contraposée pour le cas stratifié (§~4.3). Or le cas stratifié est le seul réaliste.

\textbf{La conductance d'interface.} Absente du modèle. C'est le paramètre dont la sensibilité
est \textit{a priori} la plus susceptible d'être colinéaire à celle de $\tau$, et son absence
vide de contenu le passage de $S_\tau$ à une conclusion sur l'ajustabilité.

\textbf{La conception d'expérience.} L'article constate empiriquement un optimum en épaisseur à
1~\textmu m sans l'expliquer. Le critère $d \simeq \sqrt{2\alpha\tau} = \sqrt2\,v\tau$ du §~4.3
en rend compte et se généralise.

\subsection{Absences de citation significatives}

L'article ne cite ni la littérature d'identifiabilité en thermique photothermique (Krapez,
Krapez \& Rigollet), ni aucune analyse d'incertitude en thermoréflectance, ni la borne de
Cramér--Rao, ni la matrice de Fisher. Le formalisme de Guyer--Krumhansl est absent : le modèle
reste strictement Cattaneo, le DPL est mentionné en introduction sans être employé, et la revue
de Kovács est citée sans que le terme non local soit repris. L'article n'entre donc en
concurrence ni avec la partie~I ni avec la partie~II du présent travail.

\subsection{Formulation citable}

L'article est citable sur trois points, et sur trois points seulement :

\begin{enumerate}[leftmargin=1.6em,itemsep=2pt,topsep=2pt]
\item comme énoncé du problème ouvert de la caractérisation de $\tau$ --- aucun modèle théorique
      ni procédure expérimentale n'établit sa valeur pour un matériau donné~;
\item comme source de la forme de la réponse en phase sous CV et du calcul de $S_\tau$ pour
      trois géométries~;
\item comme constat, par ses propres auteurs, que la bande requise dépasse les capacités FDTR
      actuelles pour $\tau \lesssim 3\times10^{-9}$~s.
\end{enumerate}

La formule «~forte sensibilité en FDTR~» du résumé ne peut pas être citée isolément : elle
désigne une dérivée du signal, non une incertitude de paramètre, et le corps de l'article
contient la restriction de bande qui l'accompagne.

% ================================================================ 11
\section{Matériel non consulté}

Le matériel supplémentaire n'a pas été consulté. Il contient, d'après l'annonce de l'article :
la table 1 des propriétés thermiques de tous les matériaux employés~; l'équation de la fonction
de Green associée en coordonnées cylindriques~; la dérivation des équations du système
stratifié, en particulier l'équation~(13) du coefficient de transmission.

Les points 7 et 8 du §~9 en dépendent, ainsi que la vérification numérique des valeurs
$\alpha_{SiC}$, $v = \sqrt{\alpha/\tau}$ et $\sqrt{2\alpha\tau}$ employées au §~4.3, qui sont
ici des valeurs de littérature.

Aucun résumé de ce matériel n'est donné dans la présente fiche.

% ================================================================ 12
\section{Références de l'article utiles au projet}

Extraites de la bibliographie de l'article, avec leur fonction dans celui-ci.

\smallskip
{\small
\begin{tabularx}{\linewidth}{@{}c L{0.40\linewidth} X@{}}
\toprule
\textbf{Réf.} & \textbf{Source} & \textbf{Fonction} \\
\midrule
11 & Camacho de la Rosa, Becerril, Gómez-Farfán, Esquivel-Sirvent, \textit{Energies}
\textbf{14}, 7452 (2021) &
miroirs de Bragg pour ondes thermiques~; origine des paramètres de la section III~C \\
12 & Gandolfi, Benetti, Glorieux, Giannetti, Banfi, \textit{Int. J. Heat Mass Transf.}
\textbf{143}, 118553 (2019) & relation de dispersion des ondes de température \\
13 & Camacho de la Rosa, Esquivel-Sirvent, \textit{J. Phys. Commun.} \textbf{6}, 105003 (2022) &
causalité en conduction non-Fourier~; pertinent pour la partie~II \\
19 & R. Kovács, \textit{Phys. Rep.} \textbf{1048}, 1 (2024) &
revue des équations de la chaleur au-delà de Fourier~; pertinent pour les parties~I et~II \\
31 & Gandolfi, Giannetti, Banfi, \textit{Phys. Rev. Lett.} \textbf{125}, 265901 (2020) &
cristal tempéronique dans le graphène \\
39 & Regner, Majumdar, Malen, \textit{Rev. Sci. Instrum.} \textbf{84}, 064901 (2013) &
instrumentation FDTR large bande~; origine du plafond de 200~MHz \\
40 & Kirsch, Martin, Warzoha, McLean, Windover, Takeuchi, \textit{Rev. Sci. Instrum.}
\textbf{95}, 103006 (2024) &
guide d'instrumentation FDTR~; origine de la résolution de phase de quelques millidegrés \\
45 & J.-L. Auriault, \textit{Int. J. Eng. Sci.} \textbf{101}, 45 (2016) &
CV contre Fourier thermoélastique~; régime balistique à basse température \\
50 & Y. Ezzahri, K. Joulain, \textit{J. Appl. Phys.} \textbf{112}, 083515 (2012) &
conductivité thermique dynamique des cristaux semiconducteurs~; origine de l'ordre de grandeur
de $\tau_{SiC}$ \\
51 & D. G. Cahill, \textit{Rev. Sci. Instrum.} \textbf{75}, 5119 (2004) &
matrice de transfert, pondération par le spot, définition du coefficient de sensibilité~; cité
trois fois pour trois usages distincts \\
52 & Maity, Banerjee, Mitra, \textit{Proc. SPIE} \textbf{7792}, 77920D (2010) &
dépendance en température de la réflectance des métaux \\
53 & Wang, Park, Koh, Cahill, \textit{J. Appl. Phys.} \textbf{108}, 043507 (2010) &
thermoréflectance des transducteurs métalliques~; confondant de rugosité \\
\bottomrule
\end{tabularx}}

% ================================================================ 13
\section{Résumé de la fiche en cinq énoncés}

\begin{enumerate}[leftmargin=1.7em,itemsep=4pt,topsep=2pt]
\item Article de modélisation directe, sans données ni ajustement. La substitution
      $p \mapsto p(1+\tau p)$ y est établie indépendamment et coïncide avec celle du présent
      travail.
\item La phase du demi-espace admet la forme fermée
      $\phi = -\pi/4 + \tfrac12\arctan(\omega\tau)$, dont le seuil $\omega\tau = 1$ est le
      centre exact de transition et le maximum exact de $\mathrm{d}\phi/\mathrm{d}\ln\tau$, de
      valeur $1/4$~rad soit $14{,}32^\circ$ par facteur $e$.
\item La sensibilité est définie comme $S_\tau = \mathrm{d}\phi/\mathrm{d}\ln\tau$, dérivée
      logarithmique à un seul paramètre. L'article en déduit l'ajustabilité de $\tau$~; la
      déduction n'est pas valide, faute de jacobienne complète, de corrélations, de bruit et de
      conductance d'interface.
\item Les auteurs reconnaissent dans le corps du texte que la bande requise dépasse les
      capacités FDTR courantes pour $\tau \lesssim 3\times10^{-9}$~s. Le seuil instrumental, à
      200~MHz, est $\tau_{\min} = 1/(2\pi f_{\max}) = 8{,}0\times10^{-10}$~s.
\item La conclusion de l'article revendique une méthodologie d'estimation, un ajustement et une
      quantification qui ne figurent pas dans son corps. La question de l'identifiabilité de
      $\tau$ reste entière.
\end{enumerate}

\sep

\section*{À compléter}
\addcontentsline{toc}{section}{À compléter}

\begin{itemize}[leftmargin=1.2em,itemsep=2pt,topsep=2pt]
\item Table 1 du matériel supplémentaire : propriétés thermiques employées, et reprise des
      échelles du §~4.3 avec les valeurs exactes.
\item Dérivation de la matrice de transfert du matériel supplémentaire, pour lever le point 8
      du §~9.
\item Confrontation de $S_\tau$ au calcul complet de
      $\big[(\mathbf J^{\!\top}\boldsymbol\Sigma^{-1}\mathbf J)^{-1}\big]_{\tau\tau}$ sur la même
      configuration Au/SiC/Si, conductance d'interface incluse.
\item Reprise de la figure 5(b) avec le critère $d \simeq \sqrt{2\alpha\tau}$, pour vérifier que
      l'optimum en épaisseur suit la longueur balistique lorsque $\tau$ varie.
\end{itemize}

\end{document}
